\subsection{乘积与纤维积}

设 \((\mathbf{G}_i)_{i\in I}\) 是一族带算子的群。设 \(\mathbf{G}\) 为 \(\mathbf{G}_i\) 的乘积幺半群。考虑 \(\Omega\) 在 \(\mathbf{G}\) 上的作用，定义为

\[
((x _ {i}) _ {i \in I}) ^ {\alpha} = (x _ {i} ^ {\alpha}) _ {i \in I} (\alpha \in \Omega , x _ {i} \in \mathrm{G} _ {i}).
\]

在此结构下，G 是一个带算子的群。对所有 \(i \in I\) ，投影映射 \(pr_{i}: G \to G_{i}\) 是带算子群的同态。

定义12. 上述定义的带算子群 \(\mathbf{G} = \prod_{i\in I}\mathbf{G}_i\) 是 \(\mathbf{G}_i\) 的带算子群乘积。映射 \(\mathbf{pr}_i\colon \mathbf{G}\to \mathbf{G}_i\) 称为投影同态。

带算子群乘积的一个特殊情形是群 \(G^{E}\) ，它由集合E到带算子群G的映射组成，其运算定义为：

\[
\begin{array}{r l} (f g) (x) & = f (x) g (x) \quad (f, g \in \mathbf {G} ^ {E}, x \in \mathbf {E}) \\ f ^ {\alpha} (x) & = f (x) ^ {\alpha} \quad (f \in \mathbf {G} ^ {E}, \alpha \in \Omega , x \in \mathbf {E}). \end{array}
\]

设 \((\phi_{i} \colon H \to G_{i})_{i \in I}\) 是带算子群的同态族。映射 \(h \mapsto (\phi_{i}(h))_{i \in I}\) 的 H 到 \(\prod_{i \in I} G_{i}\) 是带算子的群同态。它是唯一的同态。 \(\Phi \colon H \to \prod_{i \in I} G_{i}\) 满足 \(\mathbf{pr}_{i} \circ \Phi = \phi_{i}\) 对所有i成立。这证明了“带算子的积群”这一名称的合理性（《集合论》，IV，§2，第4节）。

设 \((\phi_i: \mathrm{H}_i \to \mathrm{G}_i)_{i \in I}\) 是带算子群的同态族。映射 \(\prod_{i \in I} \phi_i: (h_i)_{i \in I} \mapsto (\phi_i(h_i))_{i \in I}\) 于 \(\prod_{i \in I} \mathrm{H}_i\) 到 \(\prod_{i \in I} \mathrm{G}_i\) 是带算子群的同态。

命题11。设 \((\phi_i: \mathrm{H}_i \to \mathrm{G}_i)_{i \in I}\) 设为一族带算子的群同态，并令 \(\Phi = \prod_{i \in I} \phi_i\) 那么：

\begin{enumerate}
\def\labelenumi{(\alph{enumi})}
\item
  \(\operatorname{Ker}(\Phi) = \prod_{i \in I} \operatorname{Ker}(\phi_i)\) ；特别地，如果所有 \(\phi_i\) 都是单射， \(\Phi\) 是单射。
\item
  \(\operatorname{Im}(\Phi) = \prod_{i\in I}\operatorname{Im}(\phi_i)\) ；特别地，如果所有 \(\phi_i\) 都是满射， \(\Phi\) 是满射。
\end{enumerate}

这是显然的。

特别是，让 \((\mathbf{G}_{i})_{i\in\mathbf{I}}\) 设为一族带算子的群，且对所有 i，令 \(\mathbf{H}_i\) 是……的一个稳定（相应地，正规稳定）子群 \(\mathbf{G}_i\) 。乘积 \(\prod_{i\in I}\mathbf{H}_i\) 是稳定（相应地，正规稳定）子群 \(\prod_{i\in I}\mathbf{G}_i\) 以及 \(\prod_{i\in I}\mathbf{G}_i\) 到 \(\prod_{i\in I}(\mathbf{G}_i / \mathbf{H}_i)\) 定义了在过渡到商时的一个同构

\[
\left(\prod_ {i \in I} G _ {i}\right) / \left(\prod_ {i \in I} H _ {i}\right)
\]

到 \(\prod_{i\in I}(\mathrm{G}_i / \mathrm{H}_i)\) 。例如，设 J 是 I 的一个子集。子群 \(\mathbf{G}_{J}\) 于 \(\prod_{i\in I}\mathbf{G}_i\) 由 \((x_{i})_{i\in I}\) 的集合，使得 \(x_{i} = e_{i}\) 对于 \(i\notin J\) 是一个正规稳定子群。映射 \(\iota_{J}\) 它将 \(x = (x_{j})_{j\in J}\) 元素 \(y = (y_{i})_{i\in I}\) 的集合，使得 \(y_{i} = e_{i}\) 对于 \(i\notin J\) 和 \(y_{i} = x_{i}\) 对于 \(i\in J\) ，是 \(\prod_{j\in J}\mathbf{G}_j\) 到 \(\mathbf{G}_{J}\) 该映射 \(\operatorname{pr}_{I - J}\) 定义了在过渡到商群时的同构 \(\theta_J\) 商群的 \(\mathbf{G} / \mathbf{G}_J\) 到 \(\prod_{i\in I - J}\mathbf{G}_i\) 。复合 \(\operatorname{pr}_J\circ \iota_J\) 是恒等映射 \(\prod_{j\in J}\mathbf{G}_j\) . \(\mathbf{G} / \mathbf{G}_J\) 通常与 \(\prod_{i\in I - J}\mathbf{G}_i\) 等同，因为 \(\theta_J\) 和 \(\prod_{i\in J}\mathbf{G}_i\) 使得 \(\mathbf{G}_J\) 等同，因为 \(\iota_J\) .

上定义了一个处处有定义的合成律。若 \(J_{1}\) 和 \(J_{2}\) 是I的不相交子集，根据定义可知，I的每个元素 \(G_{J_{1}}\) 与……的每个元素交换 \(G_{J_{2}}\) .

定义13。设G是一个带算子的群，且 \((\mathrm{H}_{i})_{i\in I}\) G的一族正规稳定子群。设 \(p_{i}\colon G\to G/H_{i}\) 是典范同态。称G为商群族的内部积（或积） \((\mathrm{G}/\mathrm{H}_{i})\) 如果同态 \(g\mapsto(p_{i}(g))_{i\in I}\) 是 G 到 \(\prod_{i\in I}G/H_{i}\) .

设 G 和 H 是带算子的群，并设 \(\phi\) 和 \(\psi\) 是 G 到 H 的两个同态。满足 \(x\) 的 G 中的元素集合 \(\phi(x) = \psi(x)\) 是 G 的一个稳定子群，称为 \(\phi\) 和 \(\psi\) 的重合群。特别地，设 \(\phi_1: G_1 \to H\) 和 \(\phi_2: G_2 \to H\) 是带算子群的同态；这些同态 \(\phi_1 \circ \mathrm{pr}_1\) 和 \(\phi_2 \circ \mathrm{pr}_2\) 于 \(G_1 \times G_2\) 到 H 中的重合群称为 \(G_1\) 和 \(G_2\) 在 H 上相对于 \(\phi_1\) 和 \(\phi_2\) 的同态。它被记为 \(G_1 \times_H G_2\) 当没有歧义时 \(\phi_1\) 和 \(\phi_2\) 以及限制条件 \(p_1\) 和 \(p_2\) 于 \(\mathrm{pr}_1\) 和 \(\mathrm{pr}_2\) 到 \(G_1 \times_H G_2\) 也称为投影同态。 \(\phi_1 \circ p_1 = \phi_2 \circ p_2\) 。 \(G_1 \times_H G_2\) 是有序对 \((g_1, g_2) \in G_1 \times G_2\) 的集合，使得 \(\phi_1(g_1) = \phi_2(g_2)\) 的一个正规子群。若 \(f_i\) 是带算子群 K 到 \(G_i (i = 1, 2)\) 和 \(\phi_1 \circ f_1 = \phi_2 \circ f_2\) 的同态，则存在且仅存在一个同态 \(f\) 的K到 \(G_1 \times_H G_2\) 的集合，使得 \(f_i = p_i \circ f\) 对于 \(i = 1, 2\) .

